\section{Unchecked targets and the current smaller gates}

This part contains no asserted proof of the unresolved targets. A conditional reduction in the checked part means that the implication was proved; it does not mean its antecedent was established. The live question register and dated ledger give the current states, including additional deductions made after this text was first assembled.

\paragraph{U-01: general coefficient extraction and determinant-boundary ranks.}
The exact rectangular Haar transforms, ratios kernel, positive finite-group projectors and Gram-rank identities are established in R-01--R-09. The remaining general gate is a uniformly efficient positive/combinatorial realization of the multiplicity spaces and effective control of the orbit-map ranks or quotient numerator. R-27 closes the quaternary image-dimension and graded-defect gate at every size. G-20 therefore determines the eventual sign of every fixed strict four-row direction, with explicit bounds and no integer dimension transition left undecided. G-21 gives maximal permanent-image parameter dimension in every ambient variable count at least four for m at least seventeen. G-22, R-28/R-29 and G-23/G-25 now determine the eventual sign in every strict direction with $4\le r\le2m$ for $m\ge17$, with both occurrences at all sufficiently large integer stretches and no integer equality case. More than $2m$ rows, nonrectangular boundary faces beyond $r\le\min\{2n-2,m^2\}$, individual finite degrees, changing directions, and the superlinear-padding GCT regime remain unresolved. The exact eventual rectangular range starts uniformly at size fourteen; this sharp eventual threshold is not a theorem excluding every earlier finite-degree gap. G-19 supplies one unpadded discriminant copy but no total multiplicity comparison and no positive-padding separator. R-21 remains an exact fixed-ambient computation with an exponent growing with the ambient dimension. The general nonstandard quantum-group conjectures and algorithmic requirements retain their separately stated hypotheses.

\paragraph{U-02: a polynomial ROABP hitting set.}
I-30 closes the width-at-most-three, one-unknown-order case with six interpolation pairs, including arbitrarily many rank-two cuts. I-36 improves this to five pairs for the class with singular variable-dependent factors, through one-pair preservation of all nonzero maximal interval minors. I-29 proves a polynomial pair budget when there is one allowed singular rank. I-31 proves a bound controlled by the number of allowed rank levels and polynomial-size sampling for every fixed width. Its recurrence is still $2^{O(w\log w)}$ for all ranks, and sampling uses the pair count as an exponent. I-32 proves that polynomial exact coordinate-cone compression at every preserved common center is impossible. I-34 extends the prime-cocharacter repair to separated products with bounded local nullity and proves its exact one-coordinate limitation. P-14 replaces full image-cone preservation by scalar bridges for factors of rank one or full rank, reducing that case to $w-1$ pairs. General intermediate-rank matrix cancellation remains outside these divisor and product arguments. P-15 refutes a node-uniform one-pair generator already at width three with invertible factors. P-16 refutes the uncentered two-pair map on eighteen specified algebraic nodes, even for generically invertible factors, while P-14 proves that centering repairs that example. I-37 closes centered two-pair isolation for affine singular-factor width-three words. I-38 additionally permits isolated invertible affine factors. I-40 supplies a determinant-character run bound for homogeneous translation-invariant tensors and closes weights zero and one through I-37/I-38. Together with P-17 it proves the balanced two-pair theorem through nine variables, including an explicit grid. I-41 closes the weight-two class and the complete balanced problem through ten variables by a nonzero connector square coefficient and a bandwidth-three argument. I-43 closes weight three and the complete balanced theorem through eleven variables. I-44 further closes the balanced theorem through twelve variables in every characteristic and thirteen away from two. The next unrestricted small cases are thirteen variables in characteristic two and fourteen otherwise. Longer runs in unrestricted words, nonlinear singular factors, and the growing-width question remain open. The remaining target is a hitting set whose size and construction time are polynomial jointly in length, degree and growing width. Explicit joint coefficient-span preservation in I-14 remains a sufficient gate; I-03--I-04 give the stronger cone-hitting alternative. The direct Hoza transfer must respect the monomial quotient and coefficient cancellation. The independent-variable theorem, rank-profile specialization, and arithmetic encoding do not yet supply the required uniform width dependence. I-35 gives the strictly weaker rank-one-annihilator criterion and proves composition with independent seeds; the remaining bridge gate is an efficient common substitution preserving the actual restricted seed-dependent boundary functions. Arbitrary rational-boundary preservation and the stated rank-one multiplicative scalarization are refuted. The attempts register records the exact constructions tested and their failure points.

\paragraph{U-03: matrix invariants and GCT~V.}
Once U-02 is solved with the stated field and construction bounds, I-08 gives the polynomial-size separating family and finite matrix-quotient morphism. The implication is complete; the polynomial hitting-set antecedent is not. General explicit Noether normalization and general PIT have stronger hypotheses than this matrix-specific consequence. No circuit lower bound or $P$ versus $NP$ separation is asserted here.

\paragraph{U-04: a geometric or arithmetic encoding with polynomial budgets.}
The sparse and actual-sheaf constructions are proved. For arbitrary ROABPs the remaining task is to construct a sampler whose evaluation kernel avoids the bounded-width tensor variety, without requiring near-isometry of the entire coefficient norm. P-10 makes that kernel criterion exact. P-11 removes a detour: a sufficiently explicit common algebraic point is itself a hitting-set encoding, independently of GRH. P-13 constructs that point and an integer curve for every class hit by the stated centered interpolation map, with degree $(nD+1)^{2K+1}$. This gives polynomial arithmetic budgets for width three and for each fixed width through I-31. The uniform polynomial-in-width arithmetic construction remains open. A Frobenius construction would have to supply new nonvanishing structure; applying Chebotarev after assuming it does not solve the original problem.

\paragraph{U-05: unrestricted number-field random-matrix laws.}
The dependency table gives the supported statements now discharged under the stipulated inputs. Full pair correlation, spacing laws, the general arithmetic moment recipe, arbitrary automorphic target RH, and a zeta-zero quantum Hamiltonian retain their explicitly named analytic or spectral prerequisites. No literature theorem is strengthened beyond its stated support or uniformity range by changing an RH label.

\paragraph{Next work and stopping discipline.}
The attempts files record the constructions actually tried and their failure locations. The next substantive work is on the remaining joint-compression and special-fiber gates and on the restricted projected-orbit ideals. Packaging and repeated exposition are not counted as movement. Checkpoint and review times are actual elapsed times recorded in the ledger; no scheduled review is claimed before it has occurred.


\subsection{Updated gates after the generic-stabilizer and chain-anchor arguments}
R-27 closes the uniform quaternary dimension and specified Tor gate.
R-29/R-30 and G-25 give the strict-direction sign through
$2m$ rows. G-26 classifies every fixed direction through
$m$ rows. R-32/R-33 and G-27/G-28 classify every fixed
nonrectangular direction through $2m$ rows, with a further
boundary-only range $r\le\min\{2n-2,m^2\}$.
Rectangular directions beyond $m$ rows, strict directions
past $2m$ rows, changing directions and efficient finite-degree
extraction remain open. R-31 excludes the unchanged two-product
no-scaling strategy in specified $2m+1$-row directions;
it imposes no vanishing on the full cones.

P-19 proves the cubic theorem in every variable count and
characteristic. I-44 preserves all one-pole interval ranks
through five full-rank cores by its complete torus-chain
anchor calculation. The I-42 scalar step then gives weights
at most four in every characteristic and at most five away
from two. The complete balanced two-pair theorem holds
through twelve variables, and thirteen away from two.
The characteristic-two weight-five gate consists of possible
saturated five-core runs whose ranks are now preserved but
whose reset Hessian needs an additional constraint.
At longer lengths the remaining two-anchor quadratic is
explicit; both internal coefficient types can have at least
three positions. Nonlinear singular factors, general
shared-seed synchronization and uniform polynomial dependence
on growing width remain open.
